\documentclass{article}
\usepackage[utf8]{inputenc}
\usepackage[T1]{fontenc}
\usepackage{amsmath}
\usepackage{amssymb}
\usepackage{graphicx}
\usepackage{wrapfig}
\usepackage{booktabs}
\usepackage{siunitx}
\usepackage{chemfig}
\usepackage{hyperref}
\usepackage{url}

\sisetup{
  detect-all,
  per-mode=symbol,
  fraction-function=\tfrac
}

\graphicspath{{./figures/}}

\title{Physics and Chemistry: A Multidisciplinary Study of Optical Phenomena}

\author{Author Name}

\begin{document}

\maketitle

\begin{abstract}
We investigate optical phenomena in chemical systems using advanced spectroscopic
techniques. Our study reveals new insights into the interaction between light
and molecular structures.
\end{abstract}

\section{Optical Absorption}

The Beer-Lambert law describes light absorption:

\begin{wrapfigure}{r}{0.4\textwidth}
\centering
\chemfig{C(-[3]H)(-[5]H)=C(-[1]H)(-[7]H)}
\caption{Benzene ring structure.}
\label{fig:benzene}
\end{wrapfigure}

The absorbance $A$ is given by:
\begin{equation}
A(\lambda) = \epsilon(\lambda) \cdot c \cdot l \label{eq:beer_lambert}
\end{equation}

where $\epsilon$ is the molar absorptivity, $c$ is the concentration,
and $l$ is the path length.

\section{Spectroscopic Measurements}

The measured wavelengths and intensities are summarized in Table~\ref{tab:spectra}.

\begin{wraptable}{l}{0.5\textwidth}
\centering
\caption{Spectroscopic data at standard conditions.}
\label{tab:spectra}
\begin{tabular}{lS[table-format=3.2]S}
  \toprule
  Transition & \lambda (\unit\nm) & Intensity \\
  \midrule
  $\pi \to \pi^*$ & 254.0 & 1.25 \\
  $n \to \pi^*$ & 285.0 & 0.35 \\
  Charge transfer & 420.0 & 2.10 \\
  \bottomrule
\end{tabular}
\end{wraptable}

\section{Quantum Yields}

The quantum yield $\Phi$ is calculated as:
\begin{equation}
\Phi = \frac{\text{photons emitted}}{\text{photons absorbed}}
     = \frac{k_f}{k_f + k_{nr}} \label{eq:quantum_yield}
\end{equation}

The rate constants are:
\begin{align}
k_f &= A \cdot e^{-E_a/(RT)} \label{eq:k_f} \\
k_{nr} &= k_0 \cdot e^{-E_{nr}/(RT)} \label{eq:k_nr}
\end{align}

\section{Chemical Kinetics}

For the reaction \ce{A -> B -> C}, the concentration profiles follow:

\begin{equation}
\begin{cases}
[A](t) = [A]_0 e^{-k_1 t} \\
[B](t) = \dfrac{k_1[A]_0}{k_2 - k_1}\bigl(e^{-k_1 t} - e^{-k_2 t}\bigr) \\
[C](t) = [A]_0 - [A](t) - [B](t)
\end{cases} \label{eq:kinetics}
\end{equation}

\section{Conclusion}

Our findings demonstrate the applicability of spectroscopic techniques
in studying molecular structure and reaction dynamics.

\bibliographystyle{plain}
\bibliography{refs}

\end{document}
